\section{The Architect's Question}\label{the-architects-question}

\begin{quote}
How does an operator organize a fleet of complementary models---dense
and sparse, quantized and full-precision---so that every request reaches
the right model at the right cost, without over-provisioning or
under-serving?
\end{quote}

This chapter answers that question with real arithmetic, routing logic,
and operational judgment.

\section{Concept}\label{concept}

A mixed-model fleet treats each model family as a resource with distinct
cost, latency, and capability profiles. Rather than running every
request through a single 70B FP16 model on 8×H100, we layer:

\begin{itemize}
\tightlist
\item
  \textbf{Dense FP16} (baseline, high-quality)
\item
  \textbf{Dense 4-bit GGUF} (low-latency, edge-friendly)
\item
  \textbf{MoE sparse} (selective activation, higher throughput per GPU)
\item
  \textbf{Quantized variants} (4-bit/8-bit trade-offs)
\end{itemize}

The fleet operates on a \emph{routing function} R(query) → model\_id
that maps incoming requests to the cheapest model satisfying quality
constraints.

Key design principles:

\begin{enumerate}
\def\labelenumi{\arabic{enumi}.}
\tightlist
\item
  \textbf{Cost-per-token} is the primary unit of comparison across
  quantizations.
\item
  \textbf{Latency-per-token} determines interactive vs.~batch placement.
\item
  \textbf{Capability gaps} (context window, tool-use, multilingual)
  dictate minimum model tiers.
\item
  \textbf{Consolidation} saves CapEx when overlap in workload profiles
  is high; \textbf{separation} preserves isolation when workloads have
  divergent SLA or security needs.
\end{enumerate}

\section{Mental Model}\label{mental-model}

Think of the fleet as a \emph{cost-aware load balancer} with three
decision layers:

\textbf{\hyperref[tab:18.1]{Table 18-1}} --- Fleet routing decision hierarchy: capability
filter → cost--latency optimization → fallthrough.

\begin{longtable}[]{@{}
  >{\raggedright\arraybackslash}p{0.2308\linewidth}
  >{\raggedright\arraybackslash}p{0.3846\linewidth}
  >{\raggedright\arraybackslash}p{0.3846\linewidth}@{}}
\toprule\noalign{}
\begin{minipage}[b]{\linewidth}\raggedright
Layer
\end{minipage} & \begin{minipage}[b]{\linewidth}\raggedright
Question
\end{minipage} & \begin{minipage}[b]{\linewidth}\raggedright
Decision
\end{minipage} \\
\midrule\noalign{}
\endhead
\bottomrule\noalign{}
\endlastfoot
\textbf{Capability filter} & Does any model in the fleet meet the
request's feature requirements? & Eliminate models that lack the feature
(e.g., context window too small, no tool-use). \\
\textbf{Cost--latency optimization} & Among remaining models, which
minimizes weighted cost = α·token\_cost + β·latency\_cost? & Route to
the cheapest suitable model; adjust α,β per deployment mode (interactive
vs.~batch). \\
\textbf{Fallthrough} & If all suitable models are at capacity, what is
the next-best alternative? & Queue, degrade quantization, or escalate to
the next tier. \\

\label{tab:18.1}\end{longtable}

The weighting (α,β) is deployment-context dependent. A chat UI tolerates
\textasciitilde200\,ms/token; a batch ETL job tolerates
\textasciitilde2\,s/token but demands lower token cost.

\section{Worked Example}\label{worked-example}

\textbf{Scenario:} \textasciitilde2,000 users, \textasciitilde5\%
concurrent (≈100 simultaneous), \textasciitilde10 rps average /
\textasciitilde40 rps peak. Total token volume: \textasciitilde9,200 in
+ \textasciitilde300 out per hour. Base hardware: 8×H100, 640\,GB total
VRAM. Model families: a 70B dense FP16, a 70B MoE (16 experts, 2
active), and a family of 8-bit/4-bit quantized models (70B equivalent).

\subsection{3.1. Hardware capacity
mapping}\label{hardware-capacity-mapping}

\begin{longtable}[]{@{}
  >{\raggedright\arraybackslash}p{0.0938\linewidth}
  >{\raggedright\arraybackslash}p{0.2031\linewidth}
  >{\raggedright\arraybackslash}p{0.2188\linewidth}
  >{\raggedright\arraybackslash}p{0.2188\linewidth}
  >{\raggedright\arraybackslash}p{0.2656\linewidth}@{}}
\toprule\noalign{}
\begin{minipage}[b]{\linewidth}\raggedright
Model
\end{minipage} & \begin{minipage}[b]{\linewidth}\raggedright
VRAM (FP16)
\end{minipage} & \begin{minipage}[b]{\linewidth}\raggedright
VRAM (8-bit)
\end{minipage} & \begin{minipage}[b]{\linewidth}\raggedright
VRAM (4-bit)
\end{minipage} & \begin{minipage}[b]{\linewidth}\raggedright
\# fit on 8×H100
\end{minipage} \\
\midrule\noalign{}
\endhead
\bottomrule\noalign{}
\endlastfoot
70B dense FP16 & 140\,GB & N/A & N/A & 1 (weights + KV fit in the 640 GB
pool; \textasciitilde2 GPUs hold the weights) \\
70B dense 8-bit & \textasciitilde70\,GB & --- & --- & 1 (fits on 1×H100,
leaves 7 for others) \\
70B MoE FP16 (2/16 experts) & 140\,GB & N/A & N/A & 0 (same as dense) \\
70B MoE 8-bit (2 active) & \textasciitilde70\,GB & --- & --- & 1 \\
70B 4-bit GGUF & \textasciitilde35\,GB & --- & --- & \textasciitilde2
(across 2 GPUs, can batch) \\
70B 8-bit GGUF & \textasciitilde70\,GB & --- & --- & 1 \\

\label{tab:18.2}\end{longtable}

With 8×H100 (640\,GB), we can simultaneously run:

\begin{itemize}
\tightlist
\item
  \textbf{One 70B 8-bit dense} (70\,GB)
\item
  \textbf{Two 70B 4-bit GGUF} (≈70\,GB total)
\item
  \textbf{Remaining \textasciitilde500\,GB} can host smaller specialized
  models (e.g., 22B, 8B) for classification or routing.
\end{itemize}

\subsection{3.2. Per-token cost (example numbers, Q4 2026
pricing)}\label{per-token-cost-example-numbers-q4-2026-pricing}

\begin{longtable}[]{@{}
  >{\raggedright\arraybackslash}p{0.1000\linewidth}
  >{\raggedright\arraybackslash}p{0.3143\linewidth}
  >{\raggedright\arraybackslash}p{0.2857\linewidth}
  >{\raggedright\arraybackslash}p{0.3000\linewidth}@{}}
\toprule\noalign{}
\begin{minipage}[b]{\linewidth}\raggedright
Model
\end{minipage} & \begin{minipage}[b]{\linewidth}\raggedright
Token cost (USD / 1M)
\end{minipage} & \begin{minipage}[b]{\linewidth}\raggedright
Latency/token (ms)
\end{minipage} & \begin{minipage}[b]{\linewidth}\raggedright
Basis
\end{minipage} \\
\midrule\noalign{}
\endhead
\bottomrule\noalign{}
\endlastfoot
70B dense FP16 (8×H100 capital amortized) & \$2.50 & 120 & capital +
energy \\
70B dense 8-bit (1×H100) & \$1.25 & 250 & public cloud {[}1P{]} \\
70B MoE 8-bit (2 active experts) & \$1.00 & 300 & public cloud
{[}1P{]} \\
70B 4-bit GGUF (local, no GPU) & \$0.15 & 800 & energy only \\

\label{tab:18.3}\end{longtable}

\subsection{3.3. Routing arithmetic}\label{routing-arithmetic}

Each request carries its full token profile together: 9,200 input + 300
output = 9,500 tokens. \textbf{Request-level routing sends the entire
request --- input and output --- to one model.} A request cannot have
70\% of its input tokens answered by model A and the rest by model B,
unless the architecture explicitly supports model composition, cascade,
or speculative decoding (a different mechanism we do not assume here).
Under ordinary routing, we therefore split by \emph{requests}, not by
in-tokens:

For 1,000 requests/hour (≈0.28 rps --- a light illustrative load; the
arithmetic is scale-invariant in tokens):

\begin{itemize}
\tightlist
\item
  \textbf{70\% of requests} (700, ≈6.65M tokens) are routine → route to
  70B 8-bit: 6.65M × \$1.25 = \textbf{\$8.31 / hr}
\item
  \textbf{20\% of requests} (200, ≈1.9M tokens) need high-quality
  reasoning → route to 70B MoE 8-bit: 1.9M × \$1.00 = \textbf{\$1.90 /
  hr}
\item
  \textbf{10\% of requests} (100, ≈0.95M tokens) are complex /
  multilingual → route to 70B dense FP16: 0.95M × \$2.50 =
  \textbf{\$2.38 / hr}
\end{itemize}

(Each request's output tokens are generated by the same model that
processed its input. Routing output tokens to a \emph{different} model
would require an explicit cascade/composition stage --- e.g., a
generator plus a separate rewriter --- which is a distinct architecture,
not a free assumption.)

\textbf{Mixed-fleet hourly cost} = \$8.31 + \$1.90 + \$2.38 =
\textbf{\$12.59 / hour} vs.~\$23.75 / hour for single-FP16 (all 9.5M
tokens at \$2.50).

\textbf{Monthly savings} ≈ (\$23.75 − \$12.59) × 730\,hr ≈
\textbf{\$8,100 / month} (≈ 47\% reduction) with comparable quality
because the 8-bit model meets quality thresholds for \textgreater90\% of
requests.

\subsection{3.4. Latency check}\label{latency-check}

The per-request latency is that of the model serving it. The
weighted-average request latency is

\[
\bar{L} = \sum_i p_i \cdot L_i = 0.70 \times 250 + 0.20 \times 300 + 0.10 \times 120 \approx 247 \text{ ms/request}
\]

(no separate output-routing term --- output rides with its requesting
model), where \(p_i\) is the fraction of requests routed to model \(i\)
and \(L_i\) its latency. Acceptable for interactive chat (typical
threshold: ≤500 ms).

The latency and capacity must also reconcile with the canonical
workload's KV residency (\hyperref[chap:17]{Chapter 17}): routing 70\% of requests to an
8-bit model on 1×H100 changes the per-host KV ceiling accordingly, and
the fleet host count is set by the model each tier runs on. These are
{[}ILLUSTRATIVE{]} pricing and latency figures for Q4 2026; the
architect re-runs the same arithmetic with measured per-model
cost/latency from a deployment benchmark (\hyperref[chap:14]{Chapter 14}).

\subsection{3.5. Consolidation vs.~separation
decision}\label{consolidation-vs.-separation-decision}

\begin{itemize}
\tightlist
\item
  \textbf{Consolidate} when: workload profile overlap \textgreater70\%,
  latency tolerance ≥300\,ms, and cost differential \textgreater30\%.
  \emph{Verified}: our scenario consolidates 70\% of traffic to 8-bit,
  saving \textasciitilde\$8/hr.
\item
  \textbf{Keep separate} when: (a) security/classification policies
  require isolation, (b) workloads have bimodal latency needs (some
  ≤50\,ms, others batch-tolerant), or (c) model versioning cadence
  differs (e.g., frequent fine-tuning on one family, stable other).
\end{itemize}

In our example, consolidation wins because the MoE and 8-bit models
capture the same request classes at lower cost, and the 2-GPU FP16
reserve handles the tail without contention.

\section{Measurement}\label{measurement}

To operate a mixed fleet we measure four cross-cutting metrics:

\begin{longtable}[]{@{}
  >{\raggedright\arraybackslash}p{0.1795\linewidth}
  >{\raggedright\arraybackslash}p{0.4103\linewidth}
  >{\raggedright\arraybackslash}p{0.4103\linewidth}@{}}
\toprule\noalign{}
\begin{minipage}[b]{\linewidth}\raggedright
Metric
\end{minipage} & \begin{minipage}[b]{\linewidth}\raggedright
How to measure
\end{minipage} & \begin{minipage}[b]{\linewidth}\raggedright
Why it matters
\end{minipage} \\
\midrule\noalign{}
\endhead
\bottomrule\noalign{}
\endlastfoot
\textbf{Weighted average cost/token} & Σ(tokens\_i × cost\_i) /
Σ(tokens\_i) across all models & Directly tracks fleet economics; target
\textless{} \$0.002/token for competitive SaaS. \\
\textbf{Per-model throughput} & requests/sec per GPU / per CPU node &
Detects saturation; informs capacity adds vs.~routing tweaks. \\
\textbf{Routing accuracy} & \% of requests served by the \emph{intended}
target model (not fallback) & Ensures the routing function R() is
well-calibrated; low accuracy means feature gaps or cost model drift. \\
\textbf{Latency p95 per model} & p95 of token-level latency per model &
Guarante SLA per tier; p95 across the fleet = weighted p95 using routing
probabilities. \\

\label{tab:18.4}\end{longtable}

Instrumentation: add a middleware shim that logs
\texttt{model\_id,\ token\_count,\ latency\_ms,\ cost\_cents} per
request. Export to Prometheus + Grafana for real-time dashboards.

\textbf{Derived quantity example:} \emph{Weighted cost rate} = (total
monthly USD) / (total tokens processed). Start with the arithmetic
above, then update daily as actual token counts and cloud prices shift.

\section{Common Mistakes}\label{common-mistakes}

\begin{enumerate}
\def\labelenumi{\arabic{enumi}.}
\tightlist
\item
  \textbf{Treating quantization as free.} 4-bit/8-bit models reduce VRAM
  but increase CPU cycles and may lower quality on long-context tasks.
  Always measure the quality hit before routing en masse.
\item
  \textbf{One-size-fits-all routing.} A static rule (``always use
  8-bit'') fails when workload mix shifts (e.g., seasonal increase in
  multilingual queries that need FP16 precision). Use a dynamic router
  with feature-gate checks.
\item
  \textbf{Ignoring capital amortization.} FP16 on 8×H100 has high
  upfront cost; spreading that cost across a mixed fleet improves
  effective cost/token. Don't forget to annualize GPU CapEx in the
  per-token math.
\item
  \textbf{Over-provisioning the top tier.} Keeping a full FP16 GPU idle
  ``just in case'' is expensive. Keep just enough capacity for the tail
  (≈10\% of peak rps) and route the rest.
\item
  \textbf{No fallthrough plan.} If the preferred model is at capacity,
  the router must have a next‑best alternative ready; otherwise requests
  time out or queue indefinitely.
\end{enumerate}

\section{Architecture Consequence}\label{architecture-consequence}

A mixed-model fleet changes the architecture in three ways:

\begin{enumerate}
\def\labelenumi{\arabic{enumi}.}
\tightlist
\item
  \textbf{Router shim} (HTTP middleware or gRPC interceptor) sits
  between the user API and the model execution layer. It evaluates
  feature requirements, reads current capacity, and selects model\_id.
\item
  \textbf{Capacity pool abstraction} -- rather than ``GPU 0 runs model
  A,'' the system tracks ``8\,×\,H100 pool has 640\,GB VRAM; model X
  consumes 70\,GB 8-bit, model Y consumes 35\,GB 4-bit.'' The pool
  manager auto‑balances placements.
\item
  \textbf{Billing/telemetry pipeline} -- each request's cost and latency
  are tagged with model\_id, enabling downstream dashboards and alerting
  (e.g., ``cost per token rising above \$0.0025 for 3 consecutive
  hours'').
\end{enumerate}

The fleet operator becomes a \emph{cost‑latency steward} rather than a
single-model optimizer.

\section{What We Still Don't Know}\label{what-we-still-dont-know}

\begin{itemize}
\tightlist
\item
  \textbf{Quality--cost curves for new quantizations.} As 3-bit and
  sub-4-bit GGUF models emerge, the trade-off surface shifts; we need
  systematic human evals + automatic benchmarks (MT‑Bench, FLAN) across
  families.
\item
  \textbf{Optimal router policy for non‑stationary workloads.}
  Reinforcement‑learning--based routing that adapts α,β weights in real
  time is promising but underexplored in production fleets.
\item
  \textbf{Cross-model distillation opportunities.} Can a 70B 4-bit model
  be distilled from a 70B FP16 teacher without quality loss, effectively
  lowering the cost floor? Empirical work needed.
\item
  \textbf{Edge--cloud continuum effects.} When users bring their own
  4-bit‑quantized local models (via Ollama, etc.), the fleet routing
  must account for hybrid inbound/outbound token flows.
\end{itemize}

\section{End-of-Chapter Mini-Case}\label{end-of-chapter-mini-case}

\textbf{Deploy a three‑tier fleet for a customer‑support chatbot.}

\begin{itemize}
\tightlist
\item
  \textbf{Tier\,1 (4-bit GGUF):} Handles 80\% of queries (FAQ, account
  lookup). Runs on edge nodes (AMD Ryzen\,7950X, 64\,GB RAM). Cost:
  \$0.0002/token, latency \textasciitilde600\,ms.
\item
  \textbf{Tier\,2 (70B 8-bit dense):} Handles 15\% of queries (order
  status, policy look‑up). Runs on 1×H100. Cost: \$0.0015/token, latency
  \textasciitilde250\,ms.
\item
  \textbf{Tier\,3 (70B FP16):} Handles 5\% of queries (complex
  escalations, multilingual). Runs on 2 GPUs reserved. Cost:
  \$0.003/token, latency \textasciitilde120\,ms.
\end{itemize}

\textbf{Arithmetic:} 1,000 users, 20\% concurrent (200 simultaneous),
\textasciitilde5 rps average / 20 rps peak. Hourly tokens:
\textasciitilde4,500 in + 500 out = 5,000.

\begin{itemize}
\tightlist
\item
  Tier\,1 processes 4,000 tokens/hr → \$0.80
\item
  Tier\,2 processes 750 tokens/hr → \$1.13
\item
  Tier\,3 processes 250 tokens/hr → \$0.75
\item
  \textbf{Total hourly cost: \$2.68} → \textbf{\$1,945 / month}.
\end{itemize}

If we had used a single 70B FP16 model: 5,000 tokens/hr × \$0.003 =
\$15/hr → \$10,800 / month. \textbf{Savings: \$8,855 / month} (≈82\%)
with latency p95 ≈ 480\,ms (within the 500\,ms SLA).

The key enabler was the router's feature gate: queries mentioning
``multilingual'' or ``technical specification'' were auto‑escalated to
Tier\,3; the rest flowed to Tier\,1 or Tier\,2 automatically.

\begin{figure}
\centering
\pandocbounded{\includegraphics[keepaspectratio,alt={\hyperref[fig:18.1]{Fig 18.1} --- Model-routing decision tree: capability filter → cost/latency gate → fallthrough, with the five specialised families as leaves {[}ILLUSTRATIVE conceptual{]}},width=\maxwidth{\textwidth},keepaspectratio]{figures/fig-18-1801.pdf}}
\caption{\hyperref[fig:18.1]{Fig 18.1} --- Model-routing decision tree: capability filter →
cost/latency gate → fallthrough, with the five specialised families as
leaves {[}ILLUSTRATIVE conceptual{]}}

\label{fig:18.1}\end{figure}

\emph{\hyperref[fig:18.1]{Fig 18.1} --- The routing decision chain. A request first passes a
capability filter (is there a specialised model that can serve it?),
then a cost/latency gate (route = f(cost, SLO)), and otherwise falls
through to a general model. Each specialised family is a leaf the router
can dispatch to. This converts the raw cost arithmetic above into the
operational routing shape an architect operates daily.}

\begin{center}\rule{0.5\linewidth}{0.5pt}\end{center}

\emph{Evidence labels: {[}1P{]} price observations from public cloud
provider pricing pages (Q4\,2026); {[}2°{]} derived per-token cost
arithmetic; {[}VERIFY{]} latency measurements from local GGUF benchmarks
on reference hardware.}
